\chapter[2023 --- Kariko et al.]{2023: Katalin Karik\'o, Drew Weissman}
\label{ch:2023_Kariko}

\section*{Main Contribution}

\textit{Discovered nucleoside base modifications that prevent inflammatory responses to synthetic mRNA, enabling the development of effective mRNA vaccines against COVID-19.}\cite{nobel-2023,lecture-2023}

\section{Official Citation}

for their discoveries concerning nucleoside base modifications that enabled the development of effective mRNA vaccines against COVID-19

\section{Background and Historical Context}

The 2023 Nobel Prize in Physiology or Medicine recognized Katalin Karik\'o and Drew Weissman ``for their discoveries concerning nucleoside base modifications that enabled the development of effective mRNA vaccines against COVID-19.'' Their work solved a fundamental problem that had plagued messenger RNA (mRNA) therapeutics for decades---the inflammatory response triggered by synthetic mRNA---and enabled the rapid development of the mRNA vaccines that played a pivotal role in controlling the COVID-19 pandemic.

The concept of using mRNA as a therapeutic agent---introducing synthetic mRNA into cells to direct the production of specific proteins---had been a goal of molecular biology and medicine since the early 1990s. The theoretical appeal of mRNA therapeutics was compelling: unlike DNA-based therapies, mRNA does not need to enter the cell nucleus, and it cannot integrate into the host genome, reducing the risk of insertional mutagenesis. Moreover, mRNA can be rapidly designed and manufactured, making it an ideal platform for responding to emerging infectious diseases.

However, the clinical development of mRNA therapeutics was hampered by a fundamental problem: synthetic mRNA introduced into the body triggered a powerful inflammatory response that both reduced protein expression and caused significant side effects. This inflammatory response was recognized by the innate immune system, which had evolved to detect foreign RNA as a marker of viral infection. The detection of foreign RNA activates pattern recognition receptors (PRRs), including Toll-like receptors (TLR3, TLR7, and TLR8) and RIG-I-like receptors, which trigger the production of inflammatory cytokines and type I interferons. This inflammatory response was a major obstacle to the clinical use of mRNA, as it reduced the efficiency of protein expression and caused potentially dangerous side effects.

The recognition of this problem motivated many researchers to find ways to reduce the immunogenicity of synthetic mRNA while maintaining its ability to direct protein production. One approach involved modifying the nucleoside building blocks of mRNA---the chemical subunits that make up the RNA chain---to reduce immune recognition. Natural RNA contains four nucleosides: adenosine, guanosine, cytidine, and uridine. In many organisms, RNA molecules contain chemically modified nucleosides that serve various functions, including the regulation of gene expression and the evasion of immune detection. The question was whether incorporating modified nucleosides into synthetic mRNA could reduce its immunogenicity without compromising its ability to be translated into protein.

\section{The Discovery/Contribution}

The breakthrough came from the collaboration between Katalin Karik\'o, a Hungarian-born biochemist, and Drew Weissman, an American immunologist, at the University of Pennsylvania in the late 1990s and early 2000s.

\textbf{Katalin Karik\'o:} Karik\'o, born in 1955 in Szabolcs-Szatm\'ar-Bereg County, Hungary, had a lifelong passion for RNA research. She earned her Ph.D. from the University of Szeged in Hungary and conducted postdoctoral research at Temple University and the University of Pennsylvania. Karik\'o was a firm believer in the therapeutic potential of mRNA, and she devoted her career to overcoming the obstacles that prevented its clinical application.

In the mid-1990s, Karik\'o began investigating the problem of mRNA immunogenicity. She hypothesized that the inflammatory response triggered by synthetic mRNA was due to its recognition by innate immune receptors, and that modifying the nucleoside composition of mRNA could reduce this recognition. Working at the University of Pennsylvania, she conducted a series of systematic experiments testing the effects of different modified nucleosides on the immunogenicity and translational capacity of mRNA.

Karik\'o and her colleague Drew Weissman showed that the incorporation of modified nucleosides, particularly pseudouridine and N1-methylpseudouridine, into synthetic mRNA dramatically reduced its immunogenicity. These modified nucleosides reduced the activation of TLR3, TLR7, and TLR8, the pattern recognition receptors most responsible for detecting foreign RNA. As a result, mRNA containing modified nucleosides triggered much lower levels of inflammatory cytokines and interferons compared to unmodified mRNA.

Critically, Karik\'o and Weissman also showed that mRNA containing modified nucleosides was translated into protein more efficiently than unmodified mRNA. This dual benefit---reduced immunogenicity and enhanced protein expression---was the key breakthrough that enabled the clinical development of mRNA therapeutics. The improved translational efficiency was attributed to several factors, including reduced activation of innate immune pathways that suppress protein synthesis and improved binding of the modified mRNA to ribosomes.

\textbf{Drew Weissman:} Weissman, an American physician-scientist born in 1959 in Lexington, Massachusetts, brought immunological expertise to the collaboration with Karik\'o. Weissman had trained at the NIH under Anthony Fauci and had developed a deep understanding of the immune system's response to foreign antigens and pathogens. His collaboration with Karik\'o combined her expertise in RNA biochemistry with his knowledge of immune recognition mechanisms.

Weissman's contribution to the discovery included the design and interpretation of experiments testing the immunogenicity of modified mRNA in immune cells, the characterization of the innate immune pathways activated by unmodified mRNA, and the development of strategies for optimizing the design of modified mRNA for therapeutic applications. He and Karik\'o published their landmark findings in a series of papers, beginning in 2005, that established the principles for creating non-inflammatory mRNA suitable for therapeutic use.

The work of Karik\'o and Weissman laid the foundation for the development of mRNA vaccines and therapeutics. Their discovery that modified nucleosides could reduce mRNA immunogenicity while enhancing protein expression was the critical enabling technology that made the rapid development of mRNA vaccines against COVID-19 possible.

When the COVID-19 pandemic emerged in early 2020, the mRNA vaccine platform---built upon the foundational work of Karik\'o and Weissman---was uniquely positioned to respond. Moderna and BioNTech/Pfizer developed mRNA vaccines encoding the SARS-CoV-2 spike protein within weeks of the publication of the viral genome sequence. The vaccines contained mRNA with N1-methylpseudouridine modifications, directly incorporating the technology developed by Karik\'o and Weissman. Clinical trials demonstrated that the mRNA vaccines were highly effective, with efficacy rates exceeding 90\% in preventing symptomatic COVID-19 infection.

The speed with which the mRNA vaccines were developed and deployed was notable in the history of vaccinology. Traditional vaccine development typically takes 10--15 years, but the mRNA vaccines were developed, tested, and authorized for emergency use within approximately 11 months of the publication of the SARS-CoV-2 genome sequence. This remarkable achievement was made possible by decades of basic research on mRNA biology and the specific breakthroughs in mRNA modification achieved by Karik\'o and Weissman.

\section{Impact on Modern Medicine}

The impact of Karik\'o and Weissman's discoveries on modern medicine has been transformative, extending far beyond the COVID-19 vaccines.

\textbf{COVID-19 vaccines:} The mRNA vaccines developed by Moderna (mRNA-1273) and BioNTech/Pfizer (BNT162b2) have been administered billions of times worldwide and have saved an estimated millions of lives during the COVID-19 pandemic. The vaccines demonstrated the power of the mRNA platform to respond rapidly to emerging infectious diseases and validated the decades of basic research that made their development possible. The success of the COVID-19 mRNA vaccines has been one of a significant medical achievements of the twenty-first century.

\textbf{Other infectious disease vaccines:} The success of the COVID-19 mRNA vaccines has stimulated the development of mRNA vaccines against other infectious diseases, including influenza, respiratory syncytial virus (RSV), cytomegalovirus (CMV), HIV, and Zika virus. Several of these vaccines are in clinical trials, and the development of mRNA vaccines against these and other pathogens is expected to continue at a rapid pace. The mRNA platform offers several advantages for vaccine development, including rapid design and manufacturing, the ability to encode multiple antigens, and the capacity to elicit both humoral and cellular immune responses.

\textbf{Cancer immunotherapy: personalized neoantigen vaccines:} The mRNA platform is being applied to the development of personalized cancer vaccines that encode tumor-specific neoantigens---mutant proteins found only on the surface of cancer cells. By sequencing a patient's tumor, identifying neoantigens, and designing an mRNA vaccine encoding those neoantigens, it may be possible to generate a targeted immune response against the patient's specific cancer. Clinical trials of personalized mRNA cancer vaccines, conducted by Moderna and BioNTech, have shown promising results in melanoma and other cancers.

\textbf{Protein replacement therapy:} The mRNA platform can be used to direct the production of therapeutic proteins in patients with genetic diseases caused by protein deficiencies. Clinical trials of mRNA-based protein replacement therapies are underway for conditions including methylmalonic acidemia, propionic acidemia, cystic fibrosis, and ornithine transcarbamylase deficiency. The ability to deliver mRNA to specific tissues and direct the production of missing or deficient proteins represents a new approach to the treatment of genetic diseases.

\textbf{Regenerative medicine and gene editing:} mRNA can be used to deliver gene editing tools, such as CRISPR-Cas9, to cells \textit{in vivo}, enabling the correction of disease-causing mutations. The transient expression of gene editing components from mRNA reduces the risk of off-target effects compared to viral delivery systems. Clinical trials of mRNA-based gene editing therapies are underway for conditions including sickle cell disease and hereditary transthyretin amyloidosis.

\textbf{Autoimmune diseases and tolerance induction:} The mRNA platform is being explored for the induction of immune tolerance in autoimmune diseases. By delivering mRNA encoding self-antigens in a non-inflammatory context, it may be possible to reprogram the immune system to tolerate self-antigens and reduce autoimmune tissue damage. This approach is being explored for conditions including celiac disease, type 1 diabetes, and multiple sclerosis.

Katalin Karik\'o and Drew Weissman were awarded the Nobel Prize in Physiology or Medicine in 2023 ``for their discoveries concerning nucleoside base modifications that enabled the development of effective mRNA vaccines against COVID-19.'' Their work has had a significant impact on global health and has opened up new possibilities for the treatment and prevention of a wide range of diseases.

\section{Legacy}

The legacy of Karik\'o and Weissman's work extends across multiple areas of medicine and biology. Their discovery that modified nucleosides could reduce the immunogenicity of synthetic mRNA was the critical breakthrough that enabled the clinical development of mRNA therapeutics, and their work has had a transformative impact on the treatment and prevention of infectious diseases.

Katalin Karik\'o's story is also a testament to the power of perseverance in the face of adversity. Throughout her career, she faced numerous setbacks, including loss of funding, rejection of grant applications, and demotion at the University of Pennsylvania. Despite these challenges, she remained committed to her research on RNA therapeutics, driven by a deep conviction in the therapeutic potential of mRNA. Her persistence and dedication ultimately led to a discovery that has saved millions of lives.

Drew Weissman's expertise in immunology was essential to the success of the collaboration with Karik\'o. His understanding of innate immune recognition mechanisms and his ability to design rigorous experiments testing the immunogenicity of modified mRNA were critical to the discovery and its development into a clinically viable technology.

Together, Karik\'o and Weissman's work has created a new paradigm in medicine---the use of mRNA as a therapeutic platform for vaccines, protein replacement therapy, gene editing, and cancer immunotherapy. The mRNA platform offers notable flexibility, speed, and versatility, and its applications are still being explored. The legacy of their discoveries will continue to shape the future of medicine for generations to come.


\section{Modern Developments}

Karikó and Weissman's mRNA modification work has enabled the COVID-19 mRNA vaccines (Pfizer-BioNTech, Moderna) that have saved millions of lives. Understanding of nucleoside modifications has enabled development of mRNA therapeutics beyond vaccines, including cancer vaccines, protein replacement therapies, and gene editing tools. Understanding of mRNA delivery has led to lipid nanoparticle engineering for tissue targeting. mRNA vaccines for influenza, RSV, and cancer are in clinical trials. Understanding of innate immune evasion has enabled development of self-amplifying RNA vaccines.


\section{Bibliography}

\begin{thebibliography}{9}
\bibitem{nobel-2023}
Nobel Prize Outreach, ``The Nobel Prize in Physiology or Medicine 2023,''\emph{NobelPrize.org}.\\
\url{https://www.nobelprize.org/prizes/medicine/2023/summary/}

\bibitem{lecture-2023}
Katalin Karik'o, ``Nobel Lecture,''\emph{NobelPrize.org}. Winner-authored primary source; downloadable from the official Nobel Prize archive.\\
\url{https://www.nobelprize.org/prizes/medicine/2023/kariko/lecture/}
\end{thebibliography}
